% theory_api_contract.tex -- 公共 API 理论层
\section{门面 API 的能力、生命周期与结果传播理论}\label{sec:api-theory}

\begin{theorynote}
API 门面不增加电气方程；本章形式化的是能力门控、异常边界、句柄复用和独立解后审计。
具体求解模型仍由 PF/OPF/动态等专项手册负责。
\end{theorynote}

\subsection{能力门控}
令 $H$、$G$、$S$ 分别表示 HiGHS、Gurobi、SCIP 可用，则整数变量能力为
\begin{equation}
supports_{int}=H\lor G\lor S.
\end{equation}
\fld{has\_*} 是编译/链接能力，\fld{supports\_*} 是由能力组合推出的特性；二者不应被
调用方当作求解成功保证。

\subsection{安全入口与异常语义}
\texttt{safe\_solve\_power\_flow} 的调用契约是
\begin{equation}
validate\_input\land(\text{validation has Error})\Rightarrow Result<T>::error,
\end{equation}
而抛异常版本透传非法输入或求解失败。门面不把后端缺失、非收敛和 fallback 改写为成功。

\subsection{会话复用}
对网络结构签名 $\sigma$，缓存只在 $\sigma_{new}=\sigma_{cached}$ 时复用符号分析；
结构变化必须重建。数值注入和设定可以在同一结构上刷新。该策略只减少装配/符号分析成本，
不改变结果向量索引空间。

\subsection{与实现的对应}
\begin{center}\footnotesize
\begin{longtable}{@{}L{7.0cm}L{8.2cm}@{}}
\toprule 理论条目 & 实现状态\\\midrule
\endfirsthead\toprule 理论条目 & 实现状态\\\midrule\endhead
能力派生布尔式 & \implfull{src/api/hacdcpf.cpp:get_solver_capabilities}\\
安全 PF 入口 & \implfull{src/api/hacdcpf.cpp:safe_solve_power_flow}\\
符号签名缓存 & \implfull{src/api/hacdcpf.cpp:hash_system_signature}\\
OPF 解后审计 & \implfull{src/api/hacdcpf.cpp:verify_opf_result}\\
跨进程会话持久化 & \implnone\\
\bottomrule
\end{longtable}
\end{center}
